\chapter*{Preface}
\addcontentsline{toc}{chapter}{Preface}

There is a particular satisfaction in variational methods: you always know
when you are wrong.
The variational principle guarantees that any trial wavefunction yields an
energy above the true ground state.
Progress is therefore unambiguous---the energy goes down or it does not---and
the direction of error is fixed by a theorem.
What remains open is only how close you can get, and at what cost.

This thesis is a record of that question applied to the quantum many-body
problem.
We take two-dimensional harmonic quantum dots, attach a physics-informed
neural network to the Schrödinger equation, and ask what breaks when one
actually builds such a pipeline and pushes it toward accuracy.
The answer is not one thing.
It is a sequence of five distinct obstacles, each with its own character,
each illuminating something about the geometry of the problem that the
formalism alone does not reveal.

The variational principle has organized computational quantum mechanics
since the early days of the field.
Neural networks now join a long line of function approximators---Slater
determinants, Jastrow factors, backflow transformations---pressed into
the same service.
What machine learning contributes is not fundamentally new in kind, but in
scale and in optimization infrastructure: the ability to represent functions
across many orders of magnitude in input scale and to update them stably
through tens of thousands of gradient steps.
The Jastrow factor was, in a sense, already a neural wavefunction.
It simply lacked the surrounding machinery.

What this thesis discovers is that the machinery matters as much as the network.
The connections between machine learning and quantum mechanics run deep enough
that the obstacles encountered here have a dual nature: they are simultaneously
physics problems and optimization problems.
Stochastic reconfiguration---invented by physicists to optimize variational
wavefunctions---is precisely natural gradient descent on the variational
manifold.
The REINFORCE estimator from reinforcement learning is the score-function
gradient of the variational energy.
The Fisher information matrix in statistics is the quantum geometric tensor
in physics.
These were not borrowed from one field by the other; they were discovered
independently by communities that rarely read each other's papers.
Recognizing that they are the same mathematics feels less like a connection
between two fields and more like the realization that there was only ever
one field, working on the same problem from different sides.

Perhaps the most instructive moments in this work are the failures.
A subtle error in the importance-sampling code went undetected for months
because its signature---deteriorating accuracy at larger particle numbers
and lower confinement---was indistinguishable from insufficient architecture.
In $Nd$ dimensions, an error in log-density grows linearly with dimension
and is catastrophically amplified in the normalized weights; the bug was
invisible where experiments ran quickly and catastrophic where they mattered.
A Langevin refinement step, intuitively motivated, turned out to introduce
a systematic gradient bias because importance-resampling weights are valid
only when samples originate from the proposal distribution, not from a
non-equilibrium Langevin path.
A second-order optimizer degraded precisely when the landscape became
ill-conditioned enough to need it, because the Fisher-matrix estimate
collapses to noise when effective sample sizes are small.
In each case, the failure was more informative than success could have been:
it identified exactly what the method was quietly assuming and what the
problem was quietly violating.

Negative results of this kind are not a concession.
They are a method.
This thesis proceeds through three ordered decisions that every practitioner
must make when applying neural networks to the many-body Schrödinger equation:
what to parameterize, how to train, and how to optimize.
The decisions are simple to state and non-trivial to make correctly, and
the record of what happens when they are made incorrectly is, in the end,
more instructive than any table of final energies.
